\chapter{\IfLanguageName{dutch}{Stand van zaken}{State of the art}}
\label{ch:stand-van-zaken}

Dit hoofdstuk bevat de literatuurstudie van deze bachelorproef. Het bespreekt de huidige stand van zaken op het gebied van AI-gebaseerde cheats, recoilcompensatie, computer vision en moderne detectiemethoden. Daarnaast wordt onderzocht hoe artificiële intelligentie gebruikt kan worden binnen FPS-games en hoe machine learning ingezet kan worden om verdachte spelersinput te detecteren.

\section{Motivatie}

De videogame-industrie behoort vandaag tot één van de snelst groeiende entertainmentsectoren ter wereld. Volgens schattingen van Grand View Research zal de wereldwijde marktwaarde van de gamingindustrie in 2026 ongeveer 370 miljard euro bedragen en verder stijgen tot meer dan 500 miljard euro tegen 2030.

\url{https://www.grandviewresearch.com/industry-analysis/gaming-industry}

Dankzij online multiplayergames kunnen spelers wereldwijd met elkaar verbonden worden zolang zij beschikken over een stabiele internetverbinding. Hierdoor worden videogames toegankelijker voor een steeds groter publiek.

Binnen deze industrie behoren \textit{First Person Shooters} (FPS-games) al jarenlang tot de populairste genres. FPS-games combineren snelle gameplay, competitiviteit en samenwerking tussen spelers. Miljoenen spelers nemen dagelijks deel aan online competitieve matches waarin reactiesnelheid, nauwkeurigheid en spelinzicht een belangrijke rol spelen.

De competitieve aard van FPS-games brengt echter ook problemen met zich mee. Een groot aantal spelers probeert via externe hulpmiddelen een oneerlijk voordeel te verkrijgen tegenover andere spelers. Dit fenomeen staat bekend als cheaten.

Moderne cheats bestaan niet langer enkel uit eenvoudige software die rechtstreeks het spel manipuleert, maar maken steeds vaker gebruik van geavanceerde technieken zoals machine learning, computer vision en externe hardware.

Naast softwarematige cheats bestaan er ook hardware-oplossingen zoals XIM-apparaten en de Cronus Zen. Deze toestellen manipuleren gebruikersinput om functies zoals automatische recoilcompensatie of aim assist mogelijk te maken. Hierdoor ontstaat een voortdurende strijd tussen ontwikkelaars van cheats en ontwikkelaars van anti-cheatsystemen.

Om de werking van dergelijke systemen beter te begrijpen, is het belangrijk eerst inzicht te krijgen in de technische werking van FPS-games en de mechanismen waarop deze cheats gebaseerd zijn.


\section{Detectie van valsspelen via machine learning}

Traditionele anti-cheatsystemen detecteren voornamelijk bekende softwareprocessen, geheugencorruptie of verdachte bestanden. Externe AI-cheats laten echter vaak geen directe sporen achter binnen het spel zelf.

Daarom verschuift de focus van moderne detectiesystemen steeds meer naar gedragsanalyse.

Bij gedragsanalyse wordt niet gekeken naar de cheatsoftware zelf, maar naar het gedrag van de speler. Hierbij kunnen kenmerken zoals muisbewegingen, reactietijden, bewegingssnelheden en richtingsveranderingen geanalyseerd worden.

Voor deze bachelorproef wordt onderzocht hoe machine learning gebruikt kan worden om verdachte muisbewegingen te detecteren.

Hiervoor wordt een dataset opgebouwd die onder andere volgende informatie bevat:

\begin{itemize}
    \item muisverplaatsing over tijd;
    \item snelheid van de beweging;
    \item hoekveranderingen;
    \item tijdsintervallen;
    \item informatie over het al dan niet actief zijn van de cheat.
\end{itemize}

Menselijke muisbewegingen bevatten doorgaans onnauwkeurigheden en variaties in snelheid. Geautomatiseerde bewegingen zijn daarentegen vaak consistenter, sneller en mathematisch nauwkeuriger.

Door een machine-learningmodel te trainen op beide soorten input kan het systeem leren onderscheid maken tussen menselijke gameplay en geautomatiseerde aiming.

Na de trainingsfase kan het model gebruikt worden om verdachte patronen in real time te analyseren en spelersgedrag te classificeren.

\section{Vergelijking met traditionele anti-cheatsystemen}

Traditionele anti-cheatsystemen zoals Valve Anti-Cheat (VAC) en BattlEye zijn voornamelijk gericht op het detecteren van bekende cheatsignaturen, geheugenmanipulatie en verdachte processen.

Hoewel deze systemen effectief zijn tegen klassieke cheats, zijn ze minder geschikt voor externe AI-gebaseerde oplossingen die enkel videobeelden analyseren en input simuleren.

Omdat dergelijke systemen geen directe wijzigingen uitvoeren binnen het spelgeheugen, worden zij moeilijker detecteerbaar via traditionele methoden.

Hierdoor groeit de interesse in AI-gebaseerde anti-cheatsystemen die gedragsanalyse combineren met machine learning om afwijkend spelersgedrag te herkennen.

\section{Conclusie}

Uit de literatuur blijkt dat AI-gebaseerde cheats steeds geavanceerder worden door gebruik te maken van computer vision, deep learning en externe hardware.

Traditionele anti-cheatsystemen ondervinden moeilijkheden bij het detecteren van deze nieuwe generatie cheats omdat zij vaak geen directe sporen achterlaten binnen het spel zelf.

Tegelijk tonen recente onderzoeken aan dat gedragsanalyse en machine learning veelbelovende technieken vormen voor het detecteren van onnatuurlijk spelersgedrag.

Deze inzichten vormen de basis voor deze bachelorproef, waarin onderzocht wordt hoe AI-gebaseerde cheats functioneren, hoe dergelijke systemen technisch opgebouwd worden en hoe machine learning gebruikt kan worden om geautomatiseerde aiming te detecteren.